\documentclass[11pt]{article}
\usepackage[a4paper,margin=1in]{geometry}
\usepackage{amsmath,amssymb,amsthm,microtype}
\usepackage[hidelinks]{hyperref}
\usepackage[nameinlink,noabbrev]{cleveref}
\usepackage[numbers]{natbib}
\newtheorem{theorem}{Theorem}[section]
\newtheorem{proposition}[theorem]{Proposition}
\newtheorem{remark}[theorem]{Remark}
\newcommand{\Zorb}{\mathcal Z_{\rm orb}}
\newcommand{\zflip}{\zeta_{\rm flip}}
\hypersetup{pdftitle={The Unique Local Counterterm between a Henon Reflection Packet and the Full Lind Zeta},pdfauthor={Liang Wang},pdfsubject={Flip zeta and relative entropy-boundary normalization},pdfkeywords={Henon map, Lind zeta, flip system, counterterm, essential singularity}}
\title{The Unique Local Counterterm between a H\'enon\\Reflection Packet and the Full Lind Zeta}
\author{Liang Wang\\\small School of Artificial Intelligence and Automation\\\small Huazhong University of Science and Technology\\\small Wuhan 430074, P. R. China}
\date{August 16, 2026}
\begin{document}
\maketitle

\begin{abstract}
The orbit-resolved odd reflection-packet Euler product for the full H\'enon
horseshoe and the Lind zeta of the full two-shift reverse action share the
same positive entropy boundary, but they do not have the same singular
ledger.  Starting from the primary-source formula
\[
 \zeta_{\rm flip}(t)=(1-2t^2)^{-1/2}
 \exp\!\left(\frac{2t+3t^2}{1-2t^2}\right),
\]
we set $u=1-\sqrt2t$ and compare it to the packet product.  The packet
accounts for exponential pole coefficient $1/\sqrt2$, whereas the full Lind
zeta has coefficient $1/\sqrt2+3/4$ and a logarithmic branch of coefficient
$-1/2$.  We prove that
\[
 u^{1/2}e^{-3/(4u)}
 \frac{\zeta_{\rm flip}(t)}{\mathcal Z_{\rm orb}(t,1)}
\]
extends holomorphically and nonvanishingly across $u=0$ as a local branch
germ.  Among counterterms $u^\beta e^{-c/u}$, the pair
$(c,\beta)=(3/4,1/2)$ is unique.  Thus odd packet data admits an exact local
reconciliation with the source zeta but does not by itself encode the full
infinite-dihedral subgroup ledger.  Global continuation, zeros, arithmetic
semantics, and operator ownership remain open.
\end{abstract}

\section{Two objects at one boundary}

For the full $n$-shift with reverse flip, \citet[Example~4.3]{KimLeePark2003}
prove
\begin{equation}\label{eq:source-n}
 \zeta_{\sigma,\rho}(t)=
 \frac1{\sqrt{1-nt^2}}
 \exp\!\left(
 \frac{nt+(n+n^2)t^2/2}{1-nt^2}
 \right).
\end{equation}
This is the Lind zeta of the full infinite-dihedral action; its definition
sums fixed points over all finite-index subgroups.  Finite-order reversal and
matrix extensions are developed in \citet{Ryu2019}.

We specialize \eqref{eq:source-n} to $n=2$:
\begin{equation}\label{eq:lind}
 \zflip(t)=(1-2t^2)^{-1/2}
 \exp\!\left(\frac{2t+3t^2}{1-2t^2}\right).
\end{equation}
The H\'enon/full-shift conjugacy \cite{DevaneyNitecki1979} also yields a
different object: the product $\Zorb(t,1)$ over primitive odd marked
reflection words.  Its exact local theorem is
\begin{equation}\label{eq:packet}
 \log\Zorb(t,1)=
 \frac{1}{\sqrt2(1-\sqrt2t)}+G_{\rm orb}(t),
\end{equation}
with $G_{\rm orb}$ analytic near $t=1/\sqrt2$.  Equation
\eqref{eq:packet} does not assert equality with \eqref{eq:lind}.

\section{The source Lind boundary ledger}

Put
\begin{equation}\label{eq:u}
 u=1-\sqrt2t,\qquad t=\frac{1-u}{\sqrt2}.
\end{equation}
Then
\begin{align}
 1-2t^2&=u(2-u),\label{eq:denominator}\\
 2t+3t^2
 &=\sqrt2(1-u)+\frac32(1-2u+u^2).\label{eq:numerator}
\end{align}

\begin{proposition}[Full Lind local form]\label{prop:lind-local}
There is a function $G_{\rm L}$, analytic at $u=0$, such that
\begin{equation}\label{eq:lind-local}
 \log\zflip(t)=
 \frac{1/\sqrt2+3/4}{u}-\frac12\log u+G_{\rm L}(u).
\end{equation}
\end{proposition}

\begin{proof}
The exponential term in \eqref{eq:lind} is the quotient of
\eqref{eq:numerator} by \eqref{eq:denominator}.  Its residue in the coordinate
$u$ is
\[
 \frac{\sqrt2+3/2}{2}=\frac1{\sqrt2}+\frac34.
\]
Subtracting this multiple of $u^{-1}$ leaves a rational function analytic at
zero.  The prefactor contributes
$-\frac12\log u-\frac12\log(2-u)$, and the latter term is analytic at zero.
\end{proof}

\begin{remark}[Subgroup information is visible]
The coefficient $1/\sqrt2$ in \eqref{eq:lind-local} agrees exactly with the
odd packet coefficient in \eqref{eq:packet}.  The remaining $3/4$ and the
square-root branch cannot be absorbed into an analytic nonzero factor.
\end{remark}

\section{Relative counterterm and uniqueness}

Subtracting \eqref{eq:packet} from \eqref{eq:lind-local} gives
\begin{equation}\label{eq:relative-log}
 \log\frac{\zflip(t)}{\Zorb(t,1)}
 =\frac{3}{4u}-\frac12\log u+G_{\rm rel}(u),
\end{equation}
where $G_{\rm rel}$ is analytic.

\begin{theorem}[Unique local counterterm]\label{thm:counterterm}
On any fixed local branch of $u^{1/2}$,
\begin{equation}\label{eq:Crel}
 C_{\rm rel}(t):=
 u^{1/2}e^{-3/(4u)}
 \frac{\zflip(t)}{\Zorb(t,1)}
\end{equation}
extends holomorphically and nonvanishingly across $u=0$.
Moreover, among all counterterms $u^\beta e^{-c/u}$, a nonzero holomorphic
extension is possible only for
\begin{equation}\label{eq:unique}
 c=\frac34,\qquad\beta=\frac12.
\end{equation}
\end{theorem}

\begin{proof}
Taking a logarithm in \eqref{eq:Crel} cancels the two singular terms in
\eqref{eq:relative-log}, leaving $G_{\rm rel}$.  Exponentiation gives the
nonzero extension.  For a general pair $(c,\beta)$, the logarithm contains
$(3/4-c)u^{-1}+(\beta-1/2)\log u+G_{\rm rel}$.  A nonzero holomorphic germ
has neither an essential exponential nor a nonintegral/logarithmic branch,
forcing \eqref{eq:unique}.
\end{proof}

The theorem is local: the square-root branch is fixed only in a neighborhood,
and no claim is made that \eqref{eq:Crel} is a globally single-valued
determinant.

\section{Exact audit and route boundary}

The accompanying certificate performs the expansion in
$\mathbb Q(\sqrt2)$, verifies the normalized algebraic factor
$u^{1/2}/\sqrt{u(2-u)}=(2-u)^{-1/2}$, and records the analytic exponent
through order six.  An independent rational audit recovers $3/4$ and
$1/2$.  Eight tests pass in normal and optimized modes, four dependencies
are locked, and 23 claim mutations are rejected.

P71 supplies a source-native local bridge: the restricted H\'enon packet can
be reconciled with the full Lind germ only after adding precisely the missing
subgroup and branch ledgers.  It does not prove global continuation, a
Fredholm realization, rational-prime weights, a trace formula, or a
self-adjoint operator.  Route A remains exploratory and Route B is not
authorized.  The next test is whether the normalized relative Taylor
coefficients satisfy a finite-state recurrence or carry a new obstruction.

\bibliographystyle{plainnat}
\bibliography{references}
\end{document}
